% =====================================================================
% AutoQP -- UML Use Case Diagram
% Standalone TikZ document. Compiles on its own; the tikzpicture below
% can also be dropped into a larger document (e.g. as a figure in the
% SRS) if a full page is not wanted.
% =====================================================================
\documentclass[tikz, border=12pt]{standalone}

\usepackage[T1]{fontenc}
\usepackage{lmodern}

\usetikzlibrary{positioning, calc, fit, backgrounds, arrows.meta, shapes.geometric}

% --- Reusable styles -------------------------------------------------
% Keeping every visual choice in \tikzset means the diagram body below
% reads as pure layout and relationships, not styling detail.
\tikzset{
  usecase/.style = {
    ellipse, draw, thick, align=center,
    minimum width=3.3cm, minimum height=1.15cm,
    font=\small, fill=gray!5, inner sep=2pt
  },
  actor/.style = {
    minimum width=1.1cm, minimum height=1.7cm,
    path picture = {
      % Simple stick figure, drawn relative to the node's own bounding
      % box so the whole style is reusable at any position.
      \draw[thick] ([yshift=0.50cm]path picture bounding box.center) circle (0.16cm);
      \draw[thick] ([yshift=0.34cm]path picture bounding box.center) -- ++(0,-0.50);
      \draw[thick] ([yshift=0.12cm]path picture bounding box.center) ++(-0.27,0) -- ++(0.54,0);
      \draw[thick] ([yshift=-0.16cm]path picture bounding box.center) -- ++(-0.22,-0.40);
      \draw[thick] ([yshift=-0.16cm]path picture bounding box.center) -- ++(0.22,-0.40);
    }
  },
  actorlabel/.style = { font=\small, align=center },
  assoc/.style   = { thick },
  include/.style = { dashed, thick, -{Stealth[length=2.2mm]} },
  stereo/.style  = { midway, fill=white, font=\tiny, inner sep=1pt }
}

\begin{document}
\begin{tikzpicture}[node distance=1.3cm and 2.6cm]

  % ------------------------------------------------------------------
  % Use cases, laid out purely with relative positioning. The shape
  % naturally narrows at "Review Staged Questions" and again at
  % "Compile Paper to PDF" -- both are single gateways every path
  % above them has to pass through, so the layout mirrors the actual
  % pipeline rather than being an arbitrary grid.
  % ------------------------------------------------------------------
  \node[usecase] (login) {Log In};
  \node[usecase, right=2.6cm of login] (manage) {Manage Subjects\\\& Assignments};

  \node[usecase, below=of login]  (genai)   {Generate Questions\\via AI};
  \node[usecase, below=of manage] (extract) {Extract Questions\\from Input};

  \coordinate (mid1) at ($(genai)!0.5!(extract)$);
  \node[usecase, below=of mid1] (review) {Review Staged\\Questions};

  \node[usecase, below left=1.3cm and 0.5cm of review]  (search) {Search Question\\Bank};
  \node[usecase, below right=1.3cm and 0.5cm of review] (fork)   {Fork Question};

  \node[usecase, below=of search] (autoassemble) {Auto-Assemble\\Paper};
  \node[usecase, below=of fork]   (refresh)      {Refresh Paper to\\Latest Version};

  \coordinate (mid2) at ($(autoassemble)!0.5!(refresh)$);
  \node[usecase, below=of mid2] (compile) {Compile Paper\\to PDF};

  % ------------------------------------------------------------------
  % System boundary -- computed automatically from the use cases via
  % the `fit` library, then pushed onto the background layer so it
  % renders behind everything without affecting draw order below.
  % ------------------------------------------------------------------
  \begin{scope}[on background layer]
    \node[draw, thick, rounded corners, inner sep=0.9cm,
          fit=(login)(manage)(genai)(extract)(review)
              (search)(fork)(autoassemble)(refresh)(compile),
          label={[font=\bfseries]north:AutoQP}] (boundary) {};
  \end{scope}

  % ------------------------------------------------------------------
  % Actors. Each pair shares one vertical lane (admin above instructor,
  % AI Provider above Compilation Service) so both are placed with a
  % single relative shift from the other, rather than independently.
  % ------------------------------------------------------------------
  \node[actor, left=4.2cm of review] (instructor) {};
  \node[actorlabel, below=0.15cm of instructor] {Instructor};

  \node[actor, above=2.6cm of instructor] (admin) {};
  \node[actorlabel, below=0.15cm of admin] {Administrator};

  \node[actor, right=4.2cm of compile] (compilesvc) {};
  \node[actorlabel, below=0.15cm of compilesvc] {Compilation\\Service};

  \node[actor, above right=1.6cm and 2.3cm of extract] (aiprovider) {};
  \node[actorlabel, below=0.15cm of aiprovider] {AI Provider};

  % ------------------------------------------------------------------
  % Associations (plain lines -- an actor simply participates).
  % ------------------------------------------------------------------
  \draw[assoc] (instructor) -- (login);
  \draw[assoc] (instructor) -- (genai);
  \draw[assoc] (instructor) -- (extract);
  \draw[assoc] (instructor) -- (review);
  \draw[assoc] (instructor) -- (search);
  \draw[assoc] (instructor) -- (fork);
  \draw[assoc] (instructor) -- (autoassemble);
  \draw[assoc] (instructor) -- (refresh);
  \draw[assoc] (instructor) -- (compile);

  \draw[assoc] (admin) -- (login);
  \draw[assoc] (admin) -- (manage);

  \draw[assoc] (aiprovider) -- (genai);
  \draw[assoc] (aiprovider) -- (extract);

  \draw[assoc] (compilesvc) -- (compile);

  % ------------------------------------------------------------------
  % <<include>> relationships -- dashed, arrow points at the included
  % use case, per standard UML convention.
  % ------------------------------------------------------------------
  \draw[include] (genai)   -- node[stereo] {\guillemotleft include\guillemotright} (review);
  \draw[include] (extract) -- node[stereo] {\guillemotleft include\guillemotright} (review);
  \draw[include] (autoassemble) -- node[stereo] {\guillemotleft include\guillemotright} (search);

\end{tikzpicture}
\end{document}
